\section*{Results: dating the undated units}
\label{sec:results}

The framework is applied to nineteen undated units: the three documentary
source composites and their sub-strata, the five books of the Pentateuch,
the Deuteronomistic prose stratum of Jeremiah, and two poems traditionally
argued to be archaic.  Models are fitted on all 25 dated units; no target
contributes to its own prediction, and all targets are dated under an
agnostic prior.  Results are given in Table~\ref{tab:targets} as conformal
date ranges rather than point estimates, for the reasons set out in
\nameref{sec:validation}.

\subsection*{The post-exilic result}

One substantive claim survives at full strength.  The posterior predictive
probability that a unit postdates 586~BCE is $\geq 0.92$ under both model
families for the P source, the D source and its Code and Frame strata, the
JE source, and Leviticus~P; the Holiness Code reaches 0.76 in the
generative family and 0.92 in ridge.  This is a statement about which side
of a single boundary the predictive mass falls on, and it is the kind of
claim the data can carry: it requires the estimator to be approximately
unbiased, not precise.

The finding is robust in three respects that matter.  It holds under both
model families, which disagree substantially about point estimates ---
median difference 59~yr, with agreement on the single most likely period
for only 9 of 19 units.  It holds under the genre correction developed
below, which if anything moves the sources \emph{later}.  And it does not
depend on the interval width, since a doubling of the conformal half-width
would still leave the great majority of predictive mass after 586~BCE for
every source composite.

\subsection*{What cannot be claimed}

No unit in Table~\ref{tab:targets} has a 68\% interval confined to a
single conventional period.  All nineteen span between two and four.  The
P source is placed at 264~BCE (68\%: 497--31~BCE) by the generative family
and 344~BCE (518--170~BCE) by ridge; these overlap comfortably, but they
straddle the Persian--Hellenistic boundary in both cases.  Earlier versions
of this work reported the P source at 361 and 404~BCE under two models and
described the result as Persian-period.  That characterisation is not
supportable.  What is supportable is that P is post-exilic and that its
predictive distribution is concentrated in the fourth and third centuries.

We note also that D~source and Deuteronomy are the same text and must not
be read as independent confirmations, and that the sub-strata of a source
are not independent of the composite that contains them.

\subsection*{The archaic poems}

The Song of the Sea (Exod~15) and the Song of Deborah (Judg~5) are the two
earliest of all nineteen targets in both model families, and the only
targets whose predictive mass sits predominantly before 586~BCE
($P_{\mathrm{post}} = 0.40$ and $0.48$ generative, $0.48$ and $0.48$
ridge).  Neither poem was used in training or in feature screening.

This is an out-of-sample check on an independently motivated hypothesis,
and the model passes it: the two texts most widely argued on comparative
Semitic and orthographic grounds to be the oldest Hebrew poetry are the two
the model places earliest, out of nineteen candidates, in both families.
We report it as corroboration of the ordinal signal rather than as a dating
of either poem, whose intervals (887--421 and 843--377~BCE, generative)
are far too wide to adjudicate the question they are usually asked to
settle.  In particular, the present analysis cannot distinguish a genuinely
early poem from a later archaizing composition; the resistant-model
diagnostic offered for that purpose in an earlier version of this work did
not function, returning zero flagged features across all 44 units, and has
been withdrawn.

\subsection*{Genre as a confound}

The dated corpus is 17/25 prophecy against 5 narrative, while the
Pentateuchal targets are predominantly narrative and legal.  A model
trained largely on prophecy and applied to narrative is exposed to exactly
the register confound this literature has long identified.

The confound is measurable in the leave-one-out residuals.  In the
generative family, narrative texts are dated $+112.4$~yr too late on
average while prophetic texts are dated $-25.1$~yr too early (Welch
$p = 0.016$).  In ridge the effect is absent ($-24.3$ versus $+31.3$~yr,
$p = 0.234$).  The disagreement between families is itself informative:
the bias is a property of the generative inversion, not of the features.

The direction matters for the substantive conclusion.  If narrative units
are systematically dated too late by roughly 112~yr in the generative
family, then the generative estimate for the P source should be corrected
\emph{earlier}, from 264~BCE to approximately 376~BCE --- which remains
firmly post-exilic, and which moves it toward the uncorrected ridge
estimate of 344~BCE rather than away from it.  The post-exilic conclusion
is therefore robust to the genre correction, and the two families converge
once the correction is applied.  We regard this as the strongest available
evidence that the post-exilic placement is not a register artifact, and we
note that it is available only because both families were fitted and
reported.
